\documentclass[11pt]{article}

\usepackage[margin=1in]{geometry}
\usepackage{hyperref}
\usepackage[shortlabels]{enumitem}
\usepackage{booktabs}
\usepackage{longtable}
\usepackage{amsmath}
\usepackage{amssymb}
\usepackage{xcolor}
\usepackage{parskip}

\hypersetup{
  colorlinks=true,
  linkcolor=blue!60!black,
  urlcolor=blue!60!black,
  citecolor=blue!60!black
}

\title{%
  \textbf{Lumen Design Constitution}\\[0.5em]
  \large A Foundational Charter for Semantic Infrastructure Research
}
\author{Lumen Project}
\date{Version 0.4 \quad---\quad \today}

\begin{document}

\maketitle

\begin{center}
  \begin{tabular}{@{}ll@{}}
    \textbf{Version:} & v0.4 \\
    \textbf{Status:}  & Foundational Working Draft (Route~A Completion; gate for Route~A2) \\
    \textbf{Language:} & English \\
  \end{tabular}
\end{center}

\vspace{1em}

\begin{abstract}
Lumen is a research project investigating whether structured code
representations---specifically those augmented with type annotations
and behavioral contracts---affect the accuracy of large language models
on code reasoning tasks, including program understanding, bug
detection, and code transformation.  In the short term, Lumen is a
controlled empirical study comparing five progressively richer code
representations as inputs to frontier LLMs.  In the long term, Lumen
aims to develop a \emph{canonical core representation} for programs
that is independent of surface syntax, serving as semantic
infrastructure on which humans, LLMs, and automated reasoning tools
can operate with equal standing.

This document is the project's design constitution: the authoritative
record of what has been decided, what remains open, and what principles
govern future decisions.  It is intended to be a living document,
revised as experiments yield evidence and as understanding deepens.
\end{abstract}

\section*{Version 0.4 Changelog (Route A Completion)}
\addcontentsline{toc}{section}{Version 0.4 Changelog}

Version 0.4 records the completion of the 90-day empirical program
(Route~A) and fixes the project's next direction. Where v0.3 was the
pre-collection gate for the R5 confirmatory studies (T1 and T3), v0.4 is
written with the full nine-cell family in hand. The v0.3 changelog and
body below are retained unaltered as the historical gate record; where
v0.4 and older text differ, v0.4 governs.

\begin{enumerate}
  \item \textbf{Route~A is complete.} The 90-day empirical study ran to
    completion: the pre-registered nine-cell confirmatory family
    (H1--H3 $\times$ T1/T2/T3) was collected and analysed, and all
    materials were released. Route~A is now closed as an empirical
    program; its findings are the input to the next phase, not a reason
    to re-open the frozen record.

  \item \textbf{Outcome: an informative negative / mixed result.} The
    central test---structure at constant information (H1: C4 vs\
    C1\textsuperscript{+})---does not reach corrected significance on any
    of the three task families. It is reported as a clean, honest
    non-rejection, not a failure to find an effect assumed to be present.
    Per task:
    \begin{itemize}[nosep]
      \item \textbf{T1 (code understanding): a clean null.} The one
        ceiling-free task returns a genuine null under adequate effective
        $n$---the strongest form of the negative result.
      \item \textbf{T2 (bug detection): direction-positive but
        underpowered.} The H1 point estimate is positive and
        direction-consistent across model-runs, but the composite score's
        ceiling/tie collapse compresses the effective $n$, and the effect
        does not survive multiplicity correction.
      \item \textbf{T3 (code transformation): ceiling-censored.} A
        high-end solve-rate ceiling censors the measurement, so T3 is
        uninformative about H1 rather than evidence against it.
    \end{itemize}
    \emph{Rationale.} The theory companion accounts for this pattern:
    under information parity the Bayes-optimal risk is format-invariant,
    so any genuine representation effect on a bounded learner is a
    capacity phenomenon of bounded size that appears only near the
    model's frontier---exactly where a ceiling-free, bounded-regime
    measurement is required to see it.

  \item \textbf{Next direction: Route~A2.} The project proceeds to
    Route~A2, a continuation of the empirical path with two coupled
    workstreams: (a)~\textbf{bounded-regime testing}---ceiling-free T2/T3
    metrics run where the theory localizes the effect (near the capacity
    frontier), so the effective $n$ is no longer destroyed by
    saturation; and (b)~a \textbf{Semantic Coding Layer prototype}
    (\texttt{docs/next-phase-design.md}, P1), the earliest product-facing
    successor to the C4 IR. Route~A2 is the active phase gated by this
    version.
    \emph{Rationale.} The Route~A diagnosis---that ceiling censoring, not
    absence of effect, is the binding limitation---makes ``measure in the
    bounded regime'' the highest-value next step; pairing it with a thin
    prototype keeps the empirical and product tracks co-evolving.

  \item \textbf{Route~B remains deferred.} The verification-core path
    (full language, compiler backend, projectional editor) stays
    explicitly out of scope. Nothing in the Route~A result argues for
    opening it now; the canonical-core-over-surface-text bet
    (Non-Negotiable~1) remains active and is reassessed only on stronger
    contrary evidence than the present data provide.
\end{enumerate}

\vspace{0.5em}

\section*{Version 0.3 Changelog (Post-Diagnostic Revision)}
\addcontentsline{toc}{section}{Version 0.3 Changelog}

Version 0.3 absorbs the outcomes of the R1 measurement-design pass and the
R2 analysis-integrity pass (Linear milestones R0--R2) so that all
forthcoming confirmatory pre-registrations (R5: T1 and T3) derive from a
single, evidence-updated source of truth. \textbf{Pre-registrations for the
T1 and T3 confirmatory studies derive from this version (v0.3) of the
Constitution and the scorers ratified herein.} Each of the five v0.3
decision areas is recorded below, including those resolved as ``no change''.

\begin{enumerate}
  \item \textbf{Success criteria and primary endpoints.}
    \emph{Decision.} The pre-registered T2 composite (0--3) is retained
    \emph{unaltered} as the frozen record of the bug-detection study; the
    diagnosed ceiling/tie collapse (paper \S5.3) is treated as a measurement
    limitation of that run, not a reason to mutate it. For \emph{future}
    task families: the T1 primary endpoint is the continuous checklist
    fraction (\texttt{score\_t1\_checklist}, fraction of ground-truth
    properties correctly stated, $\in[0,1]$); the T3 primary endpoint is the
    continuous test-pass fraction (\texttt{score\_t3}, passed/total, $\in[0,1]$,
    executed in a sandboxed subprocess with timeout). The central comparison
    remains C4 vs C1+ (H1). For any future T2 re-run, the continuous
    location metric (\texttt{score\_t2\_continuous}) is designated a
    \emph{secondary} endpoint reported alongside---not replacing---the frozen
    composite, until a re-registered T2 study adopts it as co-primary.
    \emph{Rationale.} The R1 scorers are continuous-by-construction and were
    shown (R1-1, R1-3, R1-4, R1-5) to avoid the $\le 3$-level saturation that
    collapsed the T2 effective $n$; adopting them as the T1/T3 primaries
    prevents a recurrence while leaving the frozen T2 record intact.

  \item \textbf{Round-trip validation (C4$\to$C1+$\to$C4$'$).}
    \emph{Decision.} Round-trip recovery (R4) is adopted as a
    \emph{reported diagnostic}, not yet a hard protocol gate: the
    information-parity (C1+) claim must be accompanied by published recovery
    rates, but no minimum-recovery threshold is made a pass/fail criterion in
    v0.3. A threshold may be promoted to a requirement in a later version once
    a recovery-rate distribution across the dataset is in hand (R4-2).
    \emph{Rationale.} Setting a numeric threshold before observing the
    recovery distribution would be arbitrary; transparency (reporting the
    rates) is required now, calibration of a gate is deferred.

  \item \textbf{Frozen-script regime and amendment procedure.}
    \emph{Decision.} The R2-1 bootstrap-CI back-port
    (\texttt{analyze\_confirmatory.py --with-ci}; seed 20260427, 10{,}000
    resamples, percentile) is \textbf{ratified} as a conforming amendment.
    Going forward, any change to a frozen analysis script, scorer, or dataset
    after its freeze date \emph{must} (a) carry a dated in-file amendment note,
    (b) be accompanied by a standalone amendment document under \texttt{docs/},
    and (c) leave default (un-flagged) output byte-identical unless the
    amendment is explicitly a re-freeze. This codifies the procedure whose
    absence produced the \S17.4 gap (an analysis quantity computed by an
    uncommitted sibling script).
    \emph{Rationale.} A written, mandatory amendment procedure makes
    \S17.4-class deviations impossible to introduce silently.

  \item \textbf{Sampling regime under temperature constraints.}
    \emph{Decision.} Models that reject \texttt{temperature=0} are a
    designed-for case, not an ad-hoc exception. The protocol fixes the
    deterministic setting (\texttt{temperature=0}) as the default; when a model
    refuses it, the pre-registration must (a) record the minimum permitted
    temperature actually used, (b) hold it fixed across \emph{all} conditions
    for that model so the deterministic/stochastic asymmetry never correlates
    with the representation condition, and (c) report parse-failure rates per
    condition$\times$model (R2-2) so any temperature$\times$representation
    confound is visible. Confirmatory comparisons remain \emph{within} a model.
    \emph{Rationale.} The gpt-5.5/opus-4.7 successor run exposed this gap;
    holding temperature fixed within a model and disclosing per-cell parse
    failures keeps the asymmetry from contaminating the C4-vs-C1+ contrast.

  \item \textbf{Dataset sufficiency / pre-collection power analysis.}
    \emph{Decision.} A pre-collection power analysis is now \textbf{required}
    before any R5 confirmatory collection: state the assumed effect size, the
    target power, and the resulting minimum $n$ (functions), and justify the
    chosen dataset size against it. This closes OQ4 (the Open Questions
    section) and the Phase-2 gap disclosed in paper \S3.0.
    \emph{Rationale.} The original study set $n=30$ without a power
    calculation; the T2 tie-induced effective-$n$ collapse showed the cost of
    not budgeting power in advance.
\end{enumerate}

\noindent\emph{Consistency.} These decisions are consistent with the
Experimental Protocol Addendum section and with the committed R1/R2
artifacts (the \texttt{score\_t1},
\texttt{score\_t2\_continuous}, \texttt{score\_t3} modules and the
\texttt{analyze\_confirmatory.py --with-ci} amendment). Where this changelog
and older body text differ, the changelog governs for v0.3.

\tableofcontents
\newpage

%% ============================================================
\section{Thesis}
\label{sec:thesis}

\subsection*{Short-Term Thesis (90-Day Horizon)}

Structured code representations---ranging from raw abstract syntax trees
to contract-augmented intermediate representations---affect the accuracy
with which large language models perform code understanding, bug
detection, and code transformation tasks.  The direction and magnitude
of this effect depend on both the representation and the task type.
The immediate goal of this project is to quantify these effects through
controlled experimentation using Python as the source language and at
least two frontier LLMs as subjects.

\subsection*{Long-Term Vision}

Programs should have a \emph{canonical} representation that is
independent of any particular textual syntax.  This canonical core
representation would serve as the single source of truth, with textual
notations, diagrammatic views, and LLM-optimized serializations treated
as projections---derived, potentially lossy renderings of the underlying
semantic structure.  The long-term vision of Lumen is to develop such a
representation and the surrounding infrastructure (editors, verifiers,
projection engines) so that humans, LLMs, and formal verification tools
can access program meaning on equal terms.

These two goals are related but distinct.  The short-term empirical work
tests a necessary precondition of the long-term vision: that the form in
which code is represented materially affects how well an external agent
(here, an LLM) can reason about it.  If this precondition fails---if
representation form is irrelevant to reasoning quality---then the
long-term vision loses a significant part of its motivation, and the
project should be redirected accordingly.


%% ============================================================
\section{Scope}
\label{sec:scope}

\subsection*{In Scope (Next 90 Days)}

\begin{itemize}
  \item Design and execution of a controlled experiment comparing five
        code representation formats on three LLM-facing tasks.
  \item Construction of a Python-based pipeline that converts source
        code into each of the five representations.
  \item Assembly of an evaluation dataset of 30--50 non-recursive
        Python functions, covering varied complexity levels and
        domains.
  \item A prototype schema for a contract-augmented structured
        representation (the earliest ancestor of a Lumen core IR).
  \item A pre-registered statistical analysis with confirmatory and
        exploratory components.
  \item A draft workshop paper (4--6 pages, LaTeX).
  \item Public release of all code, data, prompts, raw outputs, and
        analysis scripts in a GitHub repository.
\end{itemize}

\subsection*{Out of Scope (Explicitly Excluded for Now)}

\begin{itemize}
  \item Implementation of an interpreter or compiler for Lumen.
  \item Formal definition of a type system, effect system, or memory
        model.
  \item Design of a projectional editor or structured editing
        environment.
  \item Compiler backend targeting (WASM, LLVM, or native).
  \item Self-hosting or bootstrapping.
  \item Concurrency or parallelism models.
  \item Semantic diff algorithms.
  \item Any claim that Lumen is a usable general-purpose programming
        language.
\end{itemize}


%% ============================================================
\section{Non-Negotiables}
\label{sec:non-negotiables}

The following principles govern the project.  They are organized into
two categories.  \emph{Absolute principles} (items~2--3) cannot be
overridden.  \emph{Long-term research bets} (items~1, 4, 5) are
strong architectural preferences that motivate the project's
direction; they are not themselves empirical claims and therefore are
not directly tested by the 90-day experiment.  However, each declares
an explicit \emph{evidence threshold} at which the project will
formally reassess the bet.  This structure resolves the tension
between falsifiability (which applies to empirical claims) and
architectural commitment (which applies to design direction).

\begin{enumerate}
  \item \textbf{Canonical core over surface text (long-term research
        bet).}
        The working assumption is that a structured core IR, not
        textual source, should be the canonical representation of a
        program.  This assumption shapes the long-term vision but is
        \emph{not} what the 90-day experiment directly tests.  The
        experiment tests a narrower precondition: whether
        representation format affects LLM reasoning quality.

        \emph{Relationship to the experiment:}  A positive result
        (structured representations help) strengthens the motivation
        for this bet.  A negative result (no effect, or text is
        superior) does not logically refute the bet---the bet may
        still hold for non-LLM reasons (human cognition, formal
        verification, tooling)---but it removes the LLM-performance
        motivation, which is the primary justification offered in
        this phase.

        \emph{Evidence threshold for reassessment:}  If the
        experiment shows that C1\textsuperscript{+} (text with
        equivalent information) equals or exceeds C4 (structured IR)
        on all three primary metrics across both LLMs, the project
        will conduct a documented reassessment of this principle,
        evaluating whether sufficient non-LLM motivations remain to
        justify continued pursuit.

  \item \textbf{Falsifiability (absolute).}
        Every empirical claim must be testable and refutable.  The
        project will not defend hypotheses against contrary evidence.
        Long-term research bets (item~1) are not empirical claims
        and are therefore not subject to direct falsification by a
        single experiment.  However, the \emph{empirical motivations}
        offered for those bets are falsifiable, and if falsified, the
        bets lose their stated justification.

  \item \textbf{Rediscoverability (absolute).}
        All design decisions, experimental results, and pivots will be
        documented with their rationale.  If this project is abandoned,
        a future researcher must be able to reconstruct the reasoning
        from the repository alone.

  \item \textbf{Staged strictness (long-term research bet).}
        Any future type system or contract mechanism must support
        progressive elaboration: zero annotations at the loose end,
        full dependent types or contracts at the strict end, with a
        continuous path between them.

        \emph{Evidence threshold:}  If formal analysis demonstrates
        that staged strictness introduces unsoundness that cannot be
        repaired without abandoning the principle, it may be replaced
        by a discrete-level system with explicit boundaries.

  \item \textbf{LLM-architecture independence (long-term research
        bet).}
        The core representation must not depend on the internal
        architecture of any particular LLM.

        \emph{Evidence threshold:}  If a specific architectural
        coupling demonstrably improves LLM reasoning accuracy by a
        large margin ($d > 1.0$) and the coupling can be isolated to
        the projection layer without contaminating the Core~IR, the
        project may adopt it as a projection-layer optimization.
\end{enumerate}


%% ============================================================
\section{What Is Canonical?}
\label{sec:canonical}

The word \emph{canonical} is central to Lumen and must be used
precisely.  This section defines what it means in this project and,
equally importantly, what it does not yet mean.

\subsection*{Current Definition (90-Day Scope)}

Within the current experimental scope, ``canonical'' refers to
\textbf{structural canonicalization}: a deterministic normal form for
program representations such that two programs that differ only in
variable naming (alpha-equivalence) produce the same canonical form.

Concretely:
\begin{itemize}
  \item Bound variables are converted to de~Bruijn indices before
        canonicalization.
  \item The canonical form of a single non-recursive function is a
        directed acyclic graph (DAG) of typed nodes.
  \item A cryptographic hash is computed over the canonical form.  Two
        structurally equivalent programs (up to alpha-equivalence)
        yield the same hash.
  \item The hash of a node is computed over: (node kind, ordered child
        hashes, type-annotation hash if present, contract hash if
        present).
\end{itemize}

\subsection*{Recursion and Cycles}

A naive DAG representation cannot express recursive functions or
mutually recursive binding groups, which produce cycles in the
dependency graph.  This is not merely an implementation detail: it
directly affects the definitions of canonicality, hashing, and
equality.  Unison, the closest precedent for content-addressed code,
handles this by introducing explicit cycle-aware hash structures for
recursive definitions.  Lumen must eventually do the same.

\textbf{90-day scope restriction (hard boundary):}  The experimental
dataset is restricted to \emph{non-recursive, single-function} items.
Any function that contains direct self-recursion, mutual recursion via
helper functions, or recursive data structures in its signature is
\textbf{excluded} at the dataset curation stage.  This restriction is
enforced by a static check in the pipeline (detecting self-references
in the AST).

This restriction is acceptable for the 90-day experiment because the
experiment tests representation \emph{format} effects, not language
completeness.  Non-recursive functions provide a sufficient and
controlled test bed for this purpose.  The restriction is acknowledged
as an external validity limitation (see Section~\ref{sec:threats}).

\textbf{Long-term handling (deferred):}  Three candidate designs exist
and are recorded for future evaluation.  None is adopted now.
\begin{enumerate}[(a)]
  \item \textbf{Explicit fixed-point wrapper (Unison-style).}
        Recursive references are wrapped in an explicit
        $\texttt{fix}$ combinator node.  The hash of a recursive
        function~$f$ is $H(\texttt{fix},\, H(f_{\text{body}}))$,
        where $f_{\text{body}}$ treats the recursive reference as a
        free variable.  This is the approach closest to Unison's
        content-addressing model.
  \item \textbf{Explicit back-edge nodes.}  Recursive references are
        represented as distinguished back-edge nodes pointing to an
        enclosing binder.  The graph remains structurally acyclic with
        annotated edges.
  \item \textbf{Strongly connected component (SCC) grouping.}
        Mutually recursive bindings are grouped into SCC nodes.  Each
        SCC is hashed as a single unit with a canonical ordering of
        its members (e.g., lexicographic on de~Bruijn-normalized
        bodies).
\end{enumerate}
The choice among these will be made when the Core~IR is extended
beyond the experimental prototype, informed by both the Unison
precedent and the specific needs of Lumen's projection and
verification layers.

\subsection*{What Canonical Does Not Yet Mean}

\textbf{Semantic canonicalization}---the property that two programs
with identical input--output behavior always produce the same canonical
form---is not claimed and not attempted in this phase.
Semantic canonicalization is undecidable in general and would require
either restricting the language or introducing approximation.  It is a
long-term research question, not a 90-day deliverable.

\subsection*{Hashing}

Hashes are computed on Layer~3 (Core IR) nodes as defined in
Section~\ref{sec:levels}.  The hash function is provisional
(SHA-256 is the default) and may be replaced.  The critical property
is \textbf{determinism}: the same canonical structure must always
produce the same hash, regardless of the serialization format used to
store or transmit it.

\subsection*{Equality}

Two program fragments are \emph{structurally equal in Lumen} if and
only if their Layer~3 canonical forms produce identical hashes.  This
implies alpha-equivalence but not semantic equivalence.

\subsection*{Projection Reversibility}

A projection from Layer~3 to a derived form is \emph{lossless} if the
original Layer~3 structure can be reconstructed exactly from the
projection alone.  At least one lossless projection (the reference
serialization) must exist at all times.  Human-readable and LLM-facing
projections may be lossy for readability or token-efficiency reasons,
but the lossless serialization must always be available as a fallback.


%% ============================================================
\section{Levels of Representation}
\label{sec:levels}

Lumen distinguishes four layers of program representation.  In the
current experimental phase, only Layers~1--3 are implemented, and
Layer~3 exists only as a prototype schema.

\begin{longtable}{@{}clp{7.5cm}@{}}
  \toprule
  \textbf{Layer} & \textbf{Name} & \textbf{Description} \\
  \midrule
  \endfirsthead
  \toprule
  \textbf{Layer} & \textbf{Name} & \textbf{Description} \\
  \midrule
  \endhead
  4 & Serialization &
      Concrete byte-level encoding of a Layer~3 structure.
      Multiple serialization formats (JSON, S-expressions, binary)
      may coexist.  Must be fully reversible to Layer~3. \\
  \addlinespace
  3 & Core IR (Canonical) &
      The canonical representation.  A graph of typed nodes with
      optional type annotations and contracts.  Acyclic for
      non-recursive definitions; cycle handling is deferred (see
      Section~\ref{sec:canonical}).  Hashes are computed at this
      layer.  This is the single source of truth in the long-term
      vision. \\
  \addlinespace
  2 & Projection &
      Derived views of Layer~3 content, rendered for a specific
      audience.  Includes human-readable text, LLM-optimized
      serialization, and diagrammatic views.  May be lossless or
      lossy.  At least one lossless projection is required. \\
  \addlinespace
  1 & Surface &
      The user interface layer: editors, IDEs, CLI tools.
      Edits at this layer are translated into operations on
      Layer~3 structures.  No reversibility requirement. \\
  \bottomrule
\end{longtable}

\subsection*{Experimental Representation Conditions}

For the 90-day experiment, five conditions are tested.
C1\textsuperscript{+} is a control condition that isolates the
structure-vs.-information confound.

\begin{longtable}{@{}clp{6.5cm}@{}}
  \toprule
  \textbf{Cond.} & \textbf{Label} & \textbf{Description} \\
  \midrule
  \endfirsthead
  \toprule
  \textbf{Cond.} & \textbf{Label} & \textbf{Description} \\
  \midrule
  \endhead
  C1 & Source Text &
      Raw Python source code, exactly as a developer would read it.
      Baseline condition. \\
  \addlinespace
  C1\textsuperscript{+} & Annotated Source Text &
      Python source code with type annotations inserted inline and
      behavioral contracts appended as structured docstring comments.
      Same information content as C4 but in textual form. \\
  \addlinespace
  C2 & Raw AST &
      The Python abstract syntax tree serialized as JSON, produced by
      the standard \texttt{ast} module.  No type information. \\
  \addlinespace
  C3 & Typed AST &
      The AST enriched with type annotations inferred by
      \texttt{mypy}. \\
  \addlinespace
  C4 & Contract-Augmented IR &
      The typed AST further enriched with behavioral contracts
      (preconditions, postconditions, invariants).  Contracts are
      generated by an LLM and reviewed by the author per the contract
      review protocol (Section~\ref{sec:protocol}).  This is the
      prototype ancestor of a Lumen Core~IR. \\
  \bottomrule
\end{longtable}

The progression from C1 to C4 is intentionally cumulative: each
condition adds information to the previous one.
C1\textsuperscript{+} branches off from C4, carrying the same
information in textual form, enabling the critical comparison
that separates structure from enrichment.


%% ============================================================
\section{Research Questions and Hypotheses}
\label{sec:rq}

\subsection*{Primary Research Question}

\textbf{RQ:}  Does the representation format of code---ranging from
raw text to contract-augmented structured IR---significantly affect the
accuracy of large language models on code understanding, bug detection,
and code transformation tasks?  In particular, when information content
is held approximately constant, does structured representation
outperform textual representation?

\subsection*{Confirmatory Hypotheses (Pre-Registered)}

These hypotheses define the confirmatory analysis family.  All are
tested at $\alpha = 0.05$ with Holm--Bonferroni correction applied
across the family.

\begin{description}[style=nextline]
  \item[H1 (Structure at constant information)]
    On at least one primary metric, C4 yields significantly higher
    accuracy than C1\textsuperscript{+}.  This is the \textbf{central
    test} of the Lumen thesis: structured representation helps beyond
    mere information enrichment.
  \item[H2 (Total enrichment)]
    On at least one primary metric, C4 yields significantly higher
    accuracy than C1.
  \item[H3 (Information enrichment in text)]
    On at least one primary metric, C1\textsuperscript{+} yields
    significantly higher accuracy than C1.
\end{description}

\subsection*{Exploratory Hypotheses (Not Pre-Registered)}

The following are examined for pattern discovery but are not part of
the confirmatory family.  Results are reported with uncorrected
$p$-values and explicitly labeled ``exploratory.''

\begin{description}[style=nextline]
  \item[E1 (Structure effect without enrichment)]
    C2 vs.\ C1: does switching from text to a raw AST (no new
    information) affect accuracy?
  \item[E2 (Type effect)]
    C3 vs.\ C2: does adding type annotations to the AST improve
    accuracy?
  \item[E3 (Contract marginal effect)]
    C4 vs.\ C3: does adding contracts to the typed AST improve
    accuracy?
  \item[E4 (Task interaction)]
    The effect of representation format varies across task types.
  \item[E5 (Model interaction)]
    The effect of representation format varies across LLM models.
\end{description}

\subsection*{Null Outcome}

If H1 is not rejected (structured representation shows no advantage
over annotated text at constant information), the project will
document this as a meaningful result.  The implication would be that
\emph{information content}, not representational structure, is the
primary driver of LLM reasoning quality.  The long-term Lumen vision
would then need to be motivated on grounds other than LLM
performance---for example, human factors, formal verification, or
tooling benefits.


%% ============================================================
\section{Experimental Framing}
\label{sec:experiment}

\subsection*{Source Language}

Python is the sole source language for this experiment.  This choice
is made for three reasons: (1)~Python dominates LLM training corpora,
establishing the strongest possible baseline for text-based
representation; (2)~the \texttt{ast} standard library provides
reliable parsing; (3)~\texttt{mypy} provides a mature type inference
engine.

\subsection*{The Five Conditions}

The experiment uses five conditions.  C1 through C4 form a cumulative
enrichment chain.  C1\textsuperscript{+} is a \textbf{required}
information-parity control, not an optional addition.

\textbf{Why C1\textsuperscript{+} is required:}  Without it, the
cumulative chain C1$\to$C2$\to$C3$\to$C4 conflates two independent
variables: \emph{representation format} (text vs.\ structured) and
\emph{information content} (base vs.\ type-enriched vs.\
contract-enriched).  Comparisons like C2 vs.\ C3 or C3 vs.\ C4,
taken alone, measure the effect of \emph{adding information within a
structured format}---they do not isolate whether structure itself
matters.  C1\textsuperscript{+} holds information content
approximately equal to C4 while retaining textual format, enabling
the comparison C1\textsuperscript{+} vs.\ C4 to isolate the pure
structure effect.  Without this comparison, the central research claim
(that representational \emph{structure} helps LLM reasoning) cannot be
supported even if C4 outperforms C1.

\begin{center}
  \textrm{C1 (text)}
  $\;\xrightarrow{+\,\text{structure}}\;$
  \textrm{C2 (AST)}
  $\;\xrightarrow{+\,\text{types}}\;$
  \textrm{C3 (typed AST)}
  $\;\xrightarrow{+\,\text{contracts}}\;$
  \textrm{C4 (IR)}
\end{center}

\begin{center}
  \textrm{C1 (text)}
  $\;\xrightarrow{+\,\text{types, contracts (inline)}}\;$
  \textrm{C1\textsuperscript{+} (annotated text)}
\end{center}

\subsection*{Key Comparisons and What They Isolate}

\begin{longtable}{@{}p{4cm}p{8.5cm}@{}}
  \toprule
  \textbf{Comparison} & \textbf{What it isolates} \\
  \midrule
  \endfirsthead
  C1 vs.\ C2 &
    Pure structure effect: does switching from text to an AST
    (without adding information) help or hurt? \\
  \addlinespace
  C2 vs.\ C3 &
    Type-annotation effect within a structured representation. \\
  \addlinespace
  C3 vs.\ C4 &
    Contract effect within a structured representation. \\
  \addlinespace
  C1 vs.\ C4 &
    Total enrichment effect (endpoint comparison). \\
  \addlinespace
  \textbf{C1\textsuperscript{+} vs.\ C4} &
    \textbf{Structure effect at constant information.}  Both carry
    type and contract information; they differ only in format.  This
    is the critical comparison for the Lumen thesis. \\
  \addlinespace
  C1 vs.\ C1\textsuperscript{+} &
    Information-enrichment effect within textual format (control). \\
  \bottomrule
\end{longtable}

\subsection*{Contract Generation and Review}

Contracts for C4 and C1\textsuperscript{+} are generated by an LLM
that is \emph{not} used as a test subject in the experiment.  The
generated contracts are then reviewed and, if necessary, corrected by
the author according to the contract review protocol defined in
Section~\ref{sec:protocol}.  The same contracts (after review) are
used in both C1\textsuperscript{+} and C4, ensuring that information
content is held constant across the two conditions.

\subsection*{Separation from the Long-Term Vision}

This experiment does not test whether the Lumen canonical Core~IR is a
good \emph{design for a programming language}.  It tests a narrower
and more tractable question: whether enriching code representations
with structural, type, and contract information improves LLM task
performance---and whether the improvement, if any, is attributable to
the \emph{format} of the representation or merely to the
\emph{additional information} it carries.  Positive results on the
C1\textsuperscript{+}-vs.-C4 comparison would provide direct evidence
for the Lumen thesis.  Negative results would redirect the project.


%% ============================================================
\section{Evaluation Rubric}
\label{sec:rubric}

Three task types are used.  Metrics are classified as
\textbf{primary} (used in confirmatory statistical tests) or
\textbf{secondary} (reported for qualitative analysis only).

\subsection*{T1: Code Understanding}

\textbf{Prompt format:}  The LLM receives a function in one of the
five representations and is asked: ``Describe what this function does,
including its inputs, outputs, and any important behavioral
properties.''

\textbf{Primary metric (auto-scorable checklist):}  Each function in
the dataset has a ground-truth \emph{property checklist}: a list of
$k$ verifiable factual properties (e.g., ``returns $-1$ when the input
list is empty,'' ``runs in $O(n \log n)$ time'').  The score for each
item is the fraction of checklist properties correctly stated in the
LLM's response ($0.0$--$1.0$).  Property matching is performed by an
automated classifier (a second LLM or keyword-match script) with
spot-check audits by the author on a random 20\% sample.

\textbf{Secondary metric (human-judged holistic score, 1--5):}
A holistic rubric captures qualities (coherence, completeness of
explanation) not fully captured by the checklist.  Scored blind:
outputs are shuffled and condition labels removed before scoring.

\begin{longtable}{@{}cp{10cm}@{}}
  \toprule
  \textbf{Score} & \textbf{Criterion} \\
  \midrule
  \endfirsthead
  5 & Correct and complete: all inputs, outputs, edge cases, and
      behavioral properties accurately described. \\
  4 & Correct with minor omissions: core behavior accurate, one
      edge case or secondary property missing. \\
  3 & Partially correct: main purpose identified, but significant
      behavioral details wrong or missing. \\
  2 & Mostly incorrect: main purpose misidentified, though some
      true details present. \\
  1 & Entirely incorrect or incoherent. \\
  \bottomrule
\end{longtable}

\textbf{Reliability protocol:}  Because the holistic score is the only
human-judged secondary metric retained for T1, its reliability must be
formally established.
\begin{itemize}[nosep]
  \item The author scores 20\% of T1 items twice, separated by at
        least 48~hours, and reports intra-rater agreement (weighted
        Cohen's $\kappa$).
  \item If $\kappa < 0.6$, the holistic score is demoted to an
        appendix-only metric and excluded from all interpretive claims.
  \item If $\kappa \geq 0.6$, the holistic score remains a secondary
        metric but is never used in confirmatory analysis.
  \item The full scoring rubric, with worked examples, is published in
        the repository appendix to enable future replication.
\end{itemize}

\subsection*{T2: Bug Detection}

\textbf{Setup:}  Each function in the dataset has a correct version and
a buggy version.  Bugs are introduced manually by the author in
controlled categories: off-by-one errors, wrong comparison operators,
missing boundary checks, incorrect variable references, and swapped
arguments.  Each bug is annotated with: (a) the exact AST node or line
range, (b) the bug category, and (c) a reference fix.

\textbf{Prompt format:}  ``This function contains exactly one bug.
Identify the location of the bug, explain what is wrong, and suggest a
fix.''

\textbf{Primary metric (composite score, 0--3, auto-scored):}
\begin{itemize}[nosep]
  \item \textbf{Location} (0 or 1): The LLM identifies the correct
        line or expression.  Automated by substring/AST-node matching
        against the ground-truth annotation.
  \item \textbf{Diagnosis} (0 or 1): The LLM correctly names the bug
        category.  Automated by keyword matching against a canonical
        category label set.
  \item \textbf{Fix} (0 or 1): The proposed fix, when applied, passes
        the function's full test suite.  Fully automated.
\end{itemize}

\textbf{Secondary metric:}  None.  T2 is fully objective.

\subsection*{T3: Code Transformation}

\textbf{Setup:}  Each item consists of a function and a
natural-language specification of a required change.  Each item
includes a pre-written test suite for the post-transformation
function.

\textbf{Prompt format:}  The LLM receives the function (in one of the
five representations) plus the change specification and is asked to
produce the modified function as executable Python.

\textbf{Primary metric (test pass rate, auto-scored):}  The
transformed function is executed against the test suite.  The score is
the fraction of tests passed ($0.0$--$1.0$).  If the LLM output does
not parse or crashes, the score is $0.0$.

\textbf{Dropped metric:}  The v0.2.1 draft included a ``minimal
modification'' metric (human-judged, binary) for T3.  This metric is
\textbf{dropped} from the main evaluation.  Its subjective cost is
high relative to its contribution to the research question.  If
informally interesting patterns are noticed during analysis, they may
be noted in the paper's discussion section, but no systematic scoring
is performed.

\subsection*{Metric Summary}

\begin{longtable}{@{}lllc@{}}
  \toprule
  \textbf{Task} & \textbf{Metric} & \textbf{Scoring} &
    \textbf{Role} \\
  \midrule
  \endfirsthead
  T1 & Property-checklist fraction & Auto + audit & Primary \\
  T1 & Holistic 1--5 score & Human, blind & Secondary \\
  T2 & Composite 0--3 & Auto & Primary \\
  T3 & Test pass rate & Auto & Primary \\
  \bottomrule
\end{longtable}


%% ============================================================
\section{MVP and Success Criteria}
\label{sec:mvp}

\subsection*{MVP Definition}

The minimum viable product for the first 90 days is:
\begin{enumerate}
  \item A conversion pipeline (Python scripts) that transforms Python
        source functions into the five representation conditions
        (C1, C1\textsuperscript{+}, C2, C3, C4).
  \item An evaluation dataset of 30--50 non-recursive Python functions
        with ground-truth annotations (property checklists for T1,
        injected bugs and reference fixes for T2, transformation
        specifications and test suites for T3).
  \item Experimental results from at least two frontier LLMs across all
        five conditions and three task types.
  \item A pre-registered statistical analysis (confirmatory +
        exploratory) with appropriate correction for multiple
        comparisons.
  \item A written report suitable for submission to a workshop venue.
  \item A public GitHub repository containing all code, data, prompts,
        raw LLM outputs, and analysis scripts.
\end{enumerate}

\subsection*{Success Criteria}

Success is defined relative to the confirmatory hypotheses
(Section~\ref{sec:rq}).

\begin{longtable}{@{}lp{9cm}@{}}
  \toprule
  \textbf{Level} & \textbf{Criterion} \\
  \midrule
  \endfirsthead
  Minimum &
    The experiment completes, all data is collected, and the
    confirmatory analysis is executed.  At least one of H1--H3 is
    rejected at the corrected significance level.  All materials
    are publicly released. \\
  \addlinespace
  Moderate &
    H1 (structure at constant information) is rejected for at least
    one task type, providing direct evidence for the Lumen thesis.
    The report characterizes which tasks benefit most from structured
    representation. \\
  \addlinespace
  Strong &
    H1 is rejected for two or more task types.  Exploratory analysis
    reveals an interpretable pattern (e.g., contracts help bug
    detection more than understanding).  The results are sufficient
    for a workshop paper submission. \\
  \bottomrule
\end{longtable}

\subsection*{Informative Negative Outcomes}

If none of H1--H3 is rejected, the project documents this as a
meaningful negative result.  Specifically: if H1 fails but H3 succeeds,
the finding is that information content matters but structure does not,
which redirects Lumen toward information-centric designs.  If all three
fail, representation format is not a first-order factor, and the
long-term vision must find alternative motivation.


%% ============================================================
\section{Relation to Prior Work}
\label{sec:prior-work}

This section positions the Lumen project relative to existing work.
Novelty claims are stated carefully and limited to what the project
can actually demonstrate.

\begin{longtable}{@{}p{3.5cm}p{4.5cm}p{4.5cm}@{}}
  \toprule
  \textbf{Area} & \textbf{Relevant Work} &
    \textbf{Lumen's Relation} \\
  \midrule
  \endfirsthead
  \toprule
  \textbf{Area} & \textbf{Relevant Work} &
    \textbf{Lumen's Relation} \\
  \midrule
  \endhead

  Content addressing &
    Unison (hash-identified definitions, immutable codebase) &
    Lumen adopts content addressing as a design principle.
    Differs in that Unison retains text as the primary authoring
    surface; Lumen demotes text to a projection. \\
  \addlinespace

  Projectional editing &
    JetBrains MPS, Intentional Programming, Barista &
    Lumen shares the idea of multiple views over a canonical
    representation.  Differs in explicitly targeting LLMs as a
    first-class consumer of projections, which prior systems
    did not consider. \\
  \addlinespace

  Multi-level IR &
    MLIR, GraalVM Truffle &
    These are compiler-infrastructure IRs optimized for
    transformation passes.  Lumen's Core~IR serves a different
    purpose: human--LLM--verifier interoperability, not compiler
    optimization. \\
  \addlinespace

  Typed holes / partial programs &
    Hazel &
    Hazel supports editing and type-checking incomplete programs.
    Lumen shares the goal of working with partial information
    (staged strictness) but extends the audience to LLMs. \\
  \addlinespace

  Structured prompting for LLMs &
    Various (chain-of-thought, structured output schemas, code
    sketches) &
    Existing work explores ad hoc structuring of LLM inputs.
    Lumen proposes a principled, language-design-grounded approach
    to structured representation, tested empirically. \\
  \addlinespace

  Code representation for LLMs &
    CodeBERT, GraphCodeBERT, TreeBERT, InCoder &
    These works embed code structure into model training.
    Lumen instead modifies the \emph{input representation}
    presented to a frozen (non-fine-tuned) LLM and measures the
    effect. \\
  \bottomrule
\end{longtable}

\subsection*{Novelty Claim (Circumscribed)}

The novel contribution of the first phase is \emph{not} the
individual components (ASTs, type annotations, contracts, or content
addressing), all of which exist.  The contribution is:
\begin{quote}
  A controlled empirical study measuring how progressively enriched
  code representations affect the accuracy of frontier LLMs on code
  reasoning tasks, with an explicit control condition
  (C1\textsuperscript{+}) that separates the effect of
  representational \emph{structure} from the effect of information
  \emph{enrichment}---motivated by and feeding into the design of a
  canonical core representation for programs.
\end{quote}
This framing connects the experiment to the broader vision without
overclaiming.


%% ============================================================
\section{Threats to Validity}
\label{sec:threats}

\begin{longtable}{@{}p{3cm}p{2.5cm}p{7cm}@{}}
  \toprule
  \textbf{Threat} & \textbf{Category} & \textbf{Mitigation} \\
  \midrule
  \endfirsthead
  \toprule
  \textbf{Threat} & \textbf{Category} & \textbf{Mitigation} \\
  \midrule
  \endhead

  Training-data bias &
    Construct &
    LLMs may perform better on text because text dominates training
    data.  Mitigated by testing multiple models, reporting per-model
    results, and acknowledging this as a standing limitation. \\
  \addlinespace

  Information confound &
    Internal &
    C3/C4 add information absent from C1/C2.  The
    C1\textsuperscript{+} condition controls for this:
    C1\textsuperscript{+} vs.\ C4 isolates structure at
    approximately constant information.  ``Approximately'' is
    acknowledged: minor formatting differences remain. \\
  \addlinespace

  Small dataset &
    External &
    30--50 functions may not generalize.  Mitigated by reporting
    complexity distributions, domain coverage, effect sizes, and
    confidence intervals alongside $p$-values. \\
  \addlinespace

  Contract quality &
    Internal &
    LLM-generated contracts may contain residual errors after review.
    Mitigated by the contract review protocol
    (Section~\ref{sec:protocol}): runtime-checked assertions,
    structured provenance logging, and a quality rubric. \\
  \addlinespace

  Evaluator bias &
    Internal &
    The sole author scores secondary metrics.  Mitigated by:
    (a)~making primary metrics auto-scored, (b)~blinding secondary
    scoring, (c)~reporting intra-rater reliability for T1 holistic
    scores, (d)~demoting unreliable metrics. \\
  \addlinespace

  Prompt sensitivity &
    Construct &
    Results may depend on prompt wording.  Mitigated by holding
    prompts constant across conditions (only the representation
    varies) and publishing all templates in full. \\
  \addlinespace

  Multiple testing &
    Statistical &
    Five conditions $\times$ three tasks $\times$ two models produce
    many possible comparisons.  Mitigated by separating confirmatory
    (3 hypotheses, Holm--Bonferroni corrected) from exploratory
    analyses (uncorrected, explicitly labeled). \\
  \addlinespace

  No recursion in dataset &
    External &
    Excluding recursive functions limits generalizability.  Accepted
    as a scope limitation for the 90-day study; documented and
    flagged for future work. \\
  \bottomrule
\end{longtable}


%% ============================================================
\section{Deferred Ambitions}
\label{sec:deferred}

The following are legitimate goals of the Lumen project that are
\textbf{intentionally excluded} from the 90-day scope.  They are
recorded here to prevent scope creep and to preserve the long-term
roadmap.

\begin{enumerate}
  \item \textbf{Dependent type system.}
        A formal type system with dependent types is a long-term goal
        for expressing rich contracts.  Deferred until the prototype
        Core~IR is stable.

  \item \textbf{Algebraic effect system.}
        Effects as a unified model for IO, errors, and concurrency.
        Deferred until an interpreter exists.

  \item \textbf{Memory management model.}
        Region-based, linear, or garbage-collected memory.  Deferred
        until a compiler backend is under development.

  \item \textbf{Concurrency model.}
        Structured concurrency, actors, or other models.  Deferred.

  \item \textbf{Projectional editor.}
        A structured editor operating directly on Core~IR.  Deferred
        until Core~IR is stable and at least one lossless projection
        is implemented.

  \item \textbf{Semantic diff.}
        Meaningful diff of Core~IR graphs (rename-resistant, etc.).
        Deferred until content addressing is implemented.

  \item \textbf{Self-hosting.}
        Writing Lumen's tools in Lumen.  Deferred indefinitely.

  \item \textbf{Semantic canonicalization.}
        Normalizing semantically equivalent but structurally different
        programs to the same form.  A deep research problem deferred
        to a later phase.

  \item \textbf{General-purpose language status.}
        Lumen as a language people write production code in.  Not a
        near-term goal.

  \item \textbf{Recursion and cycle handling in Core~IR.}
        Representing recursive and mutually recursive bindings in the
        canonical graph.  Two candidate designs are recorded in
        Section~\ref{sec:canonical}; selection is deferred until the
        Core~IR extends beyond the experimental prototype.
\end{enumerate}


%% ============================================================
\section{Ninety-Day Plan}
\label{sec:plan}

The plan assumes 7--10 hours of effort per week, totaling roughly
90--130 hours over 90 days.  It is divided into three phases.

\subsection*{Phase 1: Foundations (Days 1--25)}

\textbf{Reading:}
\begin{itemize}
  \item Survey 3--5 papers on structured prompting and code
        representation for LLMs (LLM4Code, NL4SE workshop proceedings).
  \item Unison language documentation (content addressing concepts
        only; no implementation).
  \item Python \texttt{ast} module documentation.
  \item \texttt{mypy} programmatic API for type inference result
        extraction.
\end{itemize}

\textbf{Decisions:}
\begin{itemize}
  \item Finalize the evaluation dataset selection criteria (function
        length, complexity, domain distribution; non-recursive only).
  \item Write the formal specification of each representation
        condition (C1, C1\textsuperscript{+}, C2, C3, C4) as a short
        internal document.
  \item Define the T1 property-checklist format and write checklists
        for the pilot functions.
  \item Select LLMs for the experiment (default: Claude Sonnet, GPT-4o).
  \item Select the contract-generation LLM (must differ from test
        subjects).
\end{itemize}

\textbf{Building:}
\begin{itemize}
  \item Python $\to$ AST (JSON) conversion pipeline.
  \item AST $\to$ Typed AST pipeline (injecting \texttt{mypy} results).
  \item Typed AST $\to$ Contract-Augmented IR pipeline (LLM-generated
        contracts + author review per Section~\ref{sec:protocol}).
  \item C4 $\to$ C1\textsuperscript{+} projection (inline annotations
        + docstring contracts).
  \item Pilot evaluation dataset: 10--15 functions with ground truth.
\end{itemize}

\textbf{Publishable:}
\begin{itemize}
  \item GitHub repository with pipeline code and this constitution.
\end{itemize}

\subsection*{Phase 2: Pilot Experiment and Calibration (Days 26--50)}

\textbf{Reading:}
\begin{itemize}
  \item Statistical methods for within-subjects experimental design
        (Wilcoxon signed-rank, Holm--Bonferroni, Cohen's $d$, bootstrap
        confidence intervals).
  \item 1--2 exemplar workshop papers from target venues for format
        and style reference.
\end{itemize}

\textbf{Decisions:}
\begin{itemize}
  \item Based on pilot results, calibrate the automated T1
        property-checker and T2 scoring scripts.
  \item Measure intra-rater agreement on T1 holistic scores; decide
        whether to retain or demote.
  \item Finalize dataset size (minimum 30, target 50).
  \item Freeze the analysis script (confirmatory + exploratory).
\end{itemize}

\textbf{Building:}
\begin{itemize}
  \item Run pilot experiment on 10--15 functions across all conditions
        and tasks.
  \item Build automated scoring scripts for T1 checklist, T2, and T3.
  \item Complete the full evaluation dataset (30--50 functions).
  \item Automate the LLM API call and response collection pipeline.
\end{itemize}

\textbf{Publishable:}
\begin{itemize}
  \item Experimental protocol document added to the repository.
\end{itemize}

\subsection*{Phase 3: Full Experiment, Analysis, and Writing (Days 51--90)}

\textbf{Reading:}
\begin{itemize}
  \item Target venue call for papers (for format and deadline
        alignment).
\end{itemize}

\textbf{Decisions:}
\begin{itemize}
  \item Based on full results, determine next-phase direction:
        \begin{itemize}
          \item If H1 is rejected $\to$ proceed to Core~IR formal
                design (Route~B).
          \item If H1 fails but H3 succeeds $\to$ pivot to
                information-centric design.
          \item If all confirmatory hypotheses fail $\to$ reassess the
                canonical-core-first principle per its override
                threshold.
        \end{itemize}
  \item Select submission target.
\end{itemize}

\textbf{Building:}
\begin{itemize}
  \item Execute full experiment.
  \item Run confirmatory and exploratory analyses.
  \item Generate visualizations.
  \item Draft Core~IR v0.01 schema sketch informed by experimental
        findings.
  \item Draft workshop paper (4--6 pages, LaTeX).
  \item Update this Design Constitution to v0.3.
\end{itemize}

\textbf{Publishable:}
\begin{itemize}
  \item Complete repository: code, data, prompts, raw outputs, results,
        analysis scripts.
  \item arXiv preprint or workshop submission.
\end{itemize}


%% ============================================================
\section{Repository Starting Structure}
\label{sec:repo}

\begin{verbatim}
lumen/
  README.md
  LICENSE
  docs/
    design-constitution.tex
    design-constitution.pdf
    experimental-protocol.md
  src/
    pipeline/
      text_to_ast.py
      ast_to_typed_ast.py
      typed_ast_to_ir.py
      ir_to_annotated_text.py   # C4 -> C1+
      contract_generator.py
      check_no_recursion.py     # static exclusion check
    experiment/
      run_experiment.py
      score_t1_checklist.py
      score_t1_holistic.py
      score_t2.py
      score_t3.py
      analyze_confirmatory.py
      analyze_exploratory.py
    utils/
      llm_client.py
      hash.py
  data/
    functions/
      raw/              # original Python functions
      ast/              # C2: AST JSON
      typed_ast/        # C3: typed AST JSON
      ir/               # C4: contract-augmented IR
      annotated_text/   # C1+: annotated source text
    contracts/
      raw/              # contract_raw.json per function
      reviewed/         # contract_reviewed.json per function
      diffs/            # contract_diff.json per function
    ground_truth/
      checklists/       # T1 property checklists
      bugs/             # T2 bug annotations
      tests/            # T3 test suites
  results/
    raw/                # LLM outputs
    analysis/           # statistical summaries
    figures/            # plots
  paper/
    main.tex
    figures/
\end{verbatim}


%% ============================================================
\section{Decision Log}
\label{sec:decisions}

\begin{longtable}{@{}p{2cm}p{4cm}p{6.5cm}@{}}
  \toprule
  \textbf{Version} & \textbf{Decision} & \textbf{Rationale} \\
  \midrule
  \endfirsthead
  \toprule
  \textbf{Version} & \textbf{Decision} & \textbf{Rationale} \\
  \midrule
  \endhead

  v0.1 &
    Lumen adopts Core~IR as the canonical representation; text is a
    projection. &
    Founding principle.  Distinguishes Lumen from all text-first
    language projects. \\
  \addlinespace

  v0.1 &
    First target domain is LLM output verification. &
    Aligns the short-term work with the growing need for
    machine-generated code quality assurance. \\
  \addlinespace

  v0.2 &
    90-day route is empirical (Route~A), not verification-core
    design (Route~B). &
    The author is new to PL implementation; an empirical study is
    feasible within skill set and time budget.  Route~B is deferred,
    not abandoned. \\
  \addlinespace

  v0.2 &
    Source language for experiments is Python. &
    Maximizes LLM baseline accuracy (most training data) and
    tooling availability (\texttt{ast}, \texttt{mypy}). \\
  \addlinespace

  v0.2 &
    Contracts for C4 are LLM-generated, then author-reviewed. &
    Balances workload (manual writing of 30--50 contracts is
    excessive) with quality (unreviewed contracts add noise). \\
  \addlinespace

  v0.2 &
    ``Canonical'' is restricted to structural canonicalization
    (alpha-equivalence via de~Bruijn indices). &
    Semantic canonicalization is undecidable in general and out of
    scope for 90 days. \\
  \addlinespace

  v0.2 &
    At least two LLMs will be used in the experiment. &
    Single-model results cannot distinguish representation effects
    from model-specific biases. \\
  \addlinespace

  v0.2 &
    All outputs are in English and LaTeX. &
    Required for international venue submission and for attracting
    future collaborators. \\
  \addlinespace

  v0.2.1 &
    Non-negotiable~1 reclassified as ``architectural default'' with
    explicit override threshold. &
    Review identified tension between unfalsifiable architectural
    commitment and the falsifiability principle. \\
  \addlinespace

  v0.2.1 &
    Dataset restricted to non-recursive functions; recursion/cycle
    handling deferred. &
    The DAG assumption breaks under recursion.  Two candidate designs
    are recorded but not adopted.  The experiment does not require
    recursive code. \\
  \addlinespace

  v0.2.1 &
    Fifth condition C1\textsuperscript{+} added to separate structure
    from information enrichment. &
    Review identified that cumulative C1--C4 chain conflates format
    and information content.  C1\textsuperscript{+} vs.\ C4 is now
    the central test. \\
  \addlinespace

  v0.2.1 &
    Primary metrics made auto-scorable; human-judged metrics demoted
    to secondary. &
    Review identified over-reliance on subjective judgment.
    Confirmatory tests now use only auto-scored metrics. \\
  \addlinespace

  v0.2.1 &
    Confirmatory / exploratory analysis distinction introduced with
    Holm--Bonferroni correction. &
    Review identified under-specified statistical plan and risk of
    undisciplined multiple testing. \\
  \addlinespace

  v0.2.1 &
    Contract review protocol formalized with allowed edits, rejection
    criteria, and provenance logging. &
    Review identified reproducibility risk in under-specified review
    process. \\
  \addlinespace

  v0.2.2 &
    Non-negotiable~1 further refined: reclassified from
    ``architectural default'' to ``long-term research bet'' with
    explicit relationship to experiment outcomes. &
    Second review found that the v0.2.1 override-threshold language
    still appeared to shield the principle from falsification.  New
    formulation separates the bet itself (not falsifiable by one
    experiment) from its empirical motivation (falsifiable). \\
  \addlinespace

  v0.2.2 &
    C1\textsuperscript{+} elevated from ``control condition'' to
    ``required information-parity control.'' &
    Review argued this is not optional but structurally necessary
    to support the central claim about representation format. \\
  \addlinespace

  v0.2.2 &
    Recursion/cycle section strengthened: Unison precedent
    acknowledged; static exclusion check added to pipeline;
    third candidate design (fixed-point wrapper) recorded. &
    Review identified that ignoring the Unison cycle-handling
    precedent weakens credibility of the content-addressing claim. \\
  \addlinespace

  v0.2.2 &
    T3 minimal-modification metric dropped. &
    Review judged subjective cost too high relative to diagnostic
    value for the primary research question. \\
  \addlinespace

  v0.2.2 &
    T2 (bug detection) designated as primary endpoint within the
    confirmatory family. &
    Most objective scoring and most direct relevance to the
    LLM-output-verification motivation. \\
  \addlinespace

  v0.2.2 &
    Statistical plan strengthened: primary endpoint designated,
    per-model analysis classified as exploratory, mixed-effects
    alternative noted, analysis-script freeze protocol added. &
    Review identified under-specification of primary vs.\
    exploratory analyses and risk of undisciplined pairwise
    testing. \\
  \addlinespace

  v0.4 &
    Route~A (90-day empirical study) declared complete; result recorded
    as an informative negative / mixed outcome---T1 a clean null, T2
    direction-positive but underpowered, T3 ceiling-censored. &
    The pre-registered nine-cell family was collected and analysed.  The
    binding limitation is ceiling censoring of the effective sample, not
    absence of effect, consistent with the theory companion's
    information-limit prediction. \\
  \addlinespace

  v0.4 &
    Next phase set to Route~A2: bounded-regime testing plus a Semantic
    Coding Layer prototype.  Route~B (full language, compiler, editor)
    remains deferred. &
    Ceiling-free measurement near the capacity frontier is where the
    theory localizes any real representation effect; a thin product
    prototype keeps the empirical and product tracks co-evolving, and
    Route~B has no new evidential trigger. \\
  \bottomrule
\end{longtable}


%% ============================================================
\section{Open Questions}
\label{sec:open}

\begin{enumerate}
  \item \textbf{What contract language should C4 use?}
        Options include Python-like assertions, a JSON schema of
        pre/postconditions, or a custom notation.  The choice affects
        both LLM comprehension and the future Core~IR design.  The
        same notation must be used for C1\textsuperscript{+} docstring
        contracts to maintain information parity.

  \item \textbf{Which workshop venues are viable targets?}
        Candidates include LLM4Code (ICSE), MAPS (SPLASH), NL-SE,
        and ASE Tool Demo.  Deadlines and page limits must be checked.

  \item \textbf{How should the T1 automated property-checker be
        implemented?}
        Options: keyword/regex matching, a second LLM as classifier,
        or NLI-based entailment checking.  To be decided during
        Phase~1 based on pilot feasibility.

  \item \textbf{Is the dataset large enough for the confirmatory
        family?}
        With 9 confirmatory tests and Holm--Bonferroni correction, the
        effective per-test alpha is reduced.  A pilot power analysis
        (or at minimum, an effect-size sanity check from the pilot
        data) should be conducted during Phase~2 to determine whether
        30 functions suffice or 50 are needed.

  \item \textbf{Should a mixed-effects model replace or supplement the
        Wilcoxon approach?}
        The pilot will reveal whether the simpler nonparametric
        approach is adequate or whether the repeated-measures structure
        demands a regression model.  Decision deferred to Phase~2.

  \item \textbf{How should the long-term novelty of Lumen be
        positioned relative to Hazel's recent LLM work?}
        Hazel is beginning to explore typed holes in conjunction with
        LLM assistance.  The overlap with Lumen's long-term vision
        should be tracked and the differentiation sharpened before any
        long-term vision paper is attempted.
\end{enumerate}


%% ============================================================
\section{Experimental Protocol Addendum}
\label{sec:protocol}

This section consolidates the experimental protocol into a single
reference.  It is binding for the 90-day study and supersedes any
conflicting informal descriptions elsewhere in this document.


\subsection{Control Conditions}

\subsubsection*{Condition Table (Definitive)}

\begin{longtable}{@{}clccp{4.5cm}@{}}
  \toprule
  \textbf{Cond.} & \textbf{Label} & \textbf{Format} &
    \textbf{Info Level} & \textbf{Purpose} \\
  \midrule
  \endfirsthead
  \toprule
  \textbf{Cond.} & \textbf{Label} & \textbf{Format} &
    \textbf{Info Level} & \textbf{Purpose} \\
  \midrule
  \endhead
  C1 & Source Text & Text & Base & Baseline \\
  C1\textsuperscript{+} & Annotated Text & Text & Enriched &
    Information-enrichment control \\
  C2 & Raw AST & Structured & Base &
    Pure structure effect (no extra info) \\
  C3 & Typed AST & Structured & +Types &
    Marginal type-annotation effect \\
  C4 & Contract IR & Structured & Enriched &
    Full enrichment in structured form \\
  \bottomrule
\end{longtable}

The critical pair is \textbf{C1\textsuperscript{+} vs.\ C4}: both
carry the same information (types + contracts); they differ only in
format (text vs.\ structured IR).

\subsubsection*{Information-Parity Guarantee}

C1\textsuperscript{+} is constructed from C4 by projecting its type
and contract annotations back into the Python source as:
\begin{itemize}[nosep]
  \item PEP~484-style inline type annotations, and
  \item a structured docstring block containing preconditions,
        postconditions, and invariants in a fixed template.
\end{itemize}
Because C1\textsuperscript{+} is \emph{derived from} C4's annotations,
the two conditions carry the same semantic content by construction.
Minor formatting differences (e.g., JSON key names vs.\ docstring
labels) are an acknowledged residual confound.


\subsection{Contract Review Protocol}

\subsubsection*{Generation}

Contracts are generated by a designated \emph{contract-generation LLM}
that is \textbf{not} used as a test subject.  The generation prompt
requests preconditions, postconditions, and invariants in a fixed JSON
schema.

\subsubsection*{Allowed Author Edits}

During review the author may:
\begin{enumerate}[nosep]
  \item \textbf{Correct} a factually wrong predicate (e.g., change
        \texttt{len(result) >= 0} to \texttt{len(result) == len(input)}).
  \item \textbf{Delete} a contract clause that is vacuously true,
        untestable, or redundant.
  \item \textbf{Add} a missing clause if a critical behavioral property
        is absent, limited to one added clause per function.
\end{enumerate}

The author may \textbf{not}:
\begin{itemize}[nosep]
  \item Rewrite the contract from scratch.
  \item Add domain knowledge that could not reasonably be inferred from
        the function's signature and body.
  \item Introduce information not derivable from the source code.
\end{itemize}

\subsubsection*{Rejection Criteria}

A generated contract is \textbf{rejected} (and regenerated with a
fresh prompt) if:
\begin{itemize}[nosep]
  \item More than 50\% of its clauses are factually incorrect.
  \item It misidentifies the function's return type.
  \item It contains no postconditions.
\end{itemize}
At most two regeneration attempts are allowed.  If the third attempt
also fails, the function is excluded from the dataset and the
exclusion is logged.

\subsubsection*{Provenance Logging}

For every function, the repository stores:
\begin{enumerate}[nosep]
  \item The raw LLM-generated contract (\texttt{contract\_raw.json}).
  \item The reviewed contract (\texttt{contract\_reviewed.json}).
  \item A diff file (\texttt{contract\_diff.json}) recording every
        edit with a one-line justification.
  \item The number of regeneration attempts (0, 1, or 2).
\end{enumerate}

\subsubsection*{Contract Quality Rubric (Minimal)}

After review, every contract must satisfy:
\begin{longtable}{@{}lp{8.5cm}@{}}
  \toprule
  \textbf{Criterion} & \textbf{Requirement} \\
  \midrule
  \endfirsthead
  Soundness &
    Every precondition and postcondition is factually correct for
    the function as written. \\
  \addlinespace
  Non-triviality &
    At least one postcondition is non-vacuous (i.e., it rules out at
    least one logically possible return value). \\
  \addlinespace
  Runtime-checkable &
    Every clause can be expressed as a Python assertion that either
    passes or raises \texttt{AssertionError} on concrete inputs. \\
  \addlinespace
  Tested &
    The reviewed contract is executed as runtime assertions on the
    function's test suite; all assertions pass. \\
  \bottomrule
\end{longtable}


\subsection{Scoring Protocol}

\subsubsection*{Primary Metrics (Confirmatory)}

\begin{longtable}{@{}llp{7cm}@{}}
  \toprule
  \textbf{Task} & \textbf{Metric} & \textbf{Procedure} \\
  \midrule
  \endfirsthead
  T1 &
    Property-checklist fraction &
    Automated classifier checks each ground-truth property against the
    LLM response.  Author audits a random 20\% sample.  If audit
    disagreement $> 10\%$, the classifier is recalibrated and the
    full set is re-scored. \\
  \addlinespace
  T2 &
    Composite 0--3 &
    Location, diagnosis, and fix scored automatically against
    ground-truth annotations.  Fix correctness verified by running
    patched code against the test suite. \\
  \addlinespace
  T3 &
    Test pass rate &
    LLM output is parsed as Python.  If parsing fails, score = 0.
    Otherwise, the modified function is executed against the
    post-transformation test suite.  Score = fraction of tests
    passed. \\
  \bottomrule
\end{longtable}

\subsubsection*{Secondary Metrics (Exploratory)}

\begin{longtable}{@{}llp{7cm}@{}}
  \toprule
  \textbf{Task} & \textbf{Metric} & \textbf{Procedure} \\
  \midrule
  \endfirsthead
  T1 &
    Holistic 1--5 &
    Human-judged, blind.  Intra-rater reliability reported; demoted
    if $\kappa < 0.6$.  Rubric with worked examples published in
    repository appendix. \\
  \bottomrule
\end{longtable}

\noindent T3 ``minimal modification'' was included in v0.2.1 but is
\textbf{dropped} in this revision.  Its subjective cost outweighs its
diagnostic value for the primary research question.


\subsection{Statistical Analysis Plan}

\subsubsection*{Design}

Within-subjects (repeated-measures) design: every function is
presented to every LLM in every condition.  The unit of observation is
a (function, condition, model) triple.  The design is fully crossed:
$N_{\text{functions}} \times 5_{\text{conditions}} \times
2_{\text{models}} \times 3_{\text{tasks}}$.

\subsubsection*{Primary Endpoint}

The single most important test in the experiment is:

\begin{quote}
  \textbf{H1 on T2 (bug detection composite score):}  Does C4
  yield a higher median composite score than C1\textsuperscript{+}
  on the bug detection task?
\end{quote}

T2 is designated as the primary endpoint because it has the most
objective scoring (fully automated, 0--3 composite) and the most
direct relevance to the LLM-output-verification motivation.  H1 on
T1 and T3 are also confirmatory but are secondary to this test in
interpretive weight.

\subsubsection*{Confirmatory Analysis}

Three pre-registered hypotheses (H1, H2, H3) are tested on all three
primary metrics, yielding a confirmatory family of 9 tests
(3~hypotheses $\times$ 3~primary metrics).

For each test:
\begin{enumerate}[nosep]
  \item Compute the paired difference in scores between the two
        conditions (e.g., C4 minus C1\textsuperscript{+} for H1),
        aggregated per function across models (i.e., average the
        score for model~A and model~B for each function, then test
        the paired differences across functions).
  \item Test whether the median paired difference is significantly
        different from zero using the \textbf{Wilcoxon signed-rank
        test} (chosen because scores are ordinal or bounded continuous
        and normality is not assumed).
  \item Report the test statistic, exact $p$-value, and effect size
        (rank-biserial correlation $r$).
\end{enumerate}

\textbf{Holm--Bonferroni sequential correction} is applied across
the 9-test family at $\alpha = 0.05$.

\subsubsection*{Per-Model Analysis}

If the two LLMs show qualitatively different patterns (e.g., H1
rejected for model~A but not model~B), this is reported as a
model interaction finding.  Per-model results are presented in a
supplementary table but are classified as \textbf{exploratory}; they
are not part of the confirmatory family and do not receive
multiplicity correction.

\subsubsection*{Alternative Analysis (Considered but Deferred)}

A mixed-effects model (e.g., ordinal mixed-effects regression with
function as random intercept and condition $\times$ task $\times$
model as fixed effects) would more naturally handle the
repeated-measures structure and interaction terms.  This approach is
noted as a candidate for a revised analysis if the dataset is large
enough ($N \geq 40$ functions) and if the pilot reveals that the
simpler Wilcoxon approach is underpowered.  If adopted, it will be
documented as a protocol amendment before the full experiment begins.

\subsubsection*{Exploratory Analysis}

All other comparisons (E1--E5, per-task breakdowns, per-model
interactions, condition $\times$ task interaction patterns) are
reported with uncorrected $p$-values and explicitly labeled
\textsc{exploratory}.  Effect sizes and confidence intervals are
always reported.  No causal or confirmatory claims are made from
exploratory results.

\subsubsection*{Reporting Conventions}

\begin{itemize}[nosep]
  \item All $p$-values are reported to three decimal places.
  \item Effect sizes are reported with 95\% bootstrap confidence
        intervals (10{,}000 resamples).
  \item Raw data and analysis scripts are published in the repository.
  \item The analysis script is frozen after the pilot and before the
        full experiment begins.  The freeze date is recorded in the
        repository commit history.  Any post-hoc deviations from the
        frozen script are documented as protocol amendments with
        justification.
\end{itemize}


\subsection{Handling of Recursion and Cycles}

\textbf{90-day scope:}  The evaluation dataset includes only
non-recursive, single-function items.  Functions that call themselves
(directly or via mutual recursion) are excluded at the dataset curation
stage.

\textbf{Rationale:}  The experiment tests the effect of
\emph{representation format}, not language completeness.
Non-recursive functions provide a controlled test bed while avoiding
the unresolved question of how to canonicalize recursive definitions
(see Section~\ref{sec:canonical}).

\textbf{Scope limitation:}  This restriction means the results do not
generalize to recursive code.  This is documented as an external
validity limitation and flagged for future work.


\subsection{What Is Explicitly Out of Scope for This Protocol}

The following are \textbf{not} part of the 90-day experimental
protocol:

\begin{itemize}[nosep]
  \item Multi-file or multi-function programs.
  \item Class-level or module-level analysis.
  \item Recursive or mutually recursive functions.
  \item Fine-tuning or training LLMs on structured representations.
  \item Measuring LLM latency, token usage, or cost.
  \item Evaluating human (non-LLM) performance on the same tasks.
  \item Any claim about Lumen as a usable programming language.
  \item Formal verification or proof-carrying code.
\end{itemize}

These exclusions are intentional scope boundaries, not oversights.
Each is a candidate for future work.


%% ============================================================
\section*{Glossary}
\label{sec:glossary}
\addcontentsline{toc}{section}{Glossary}

\begin{description}[style=nextline, leftmargin=3.5cm]
  \item[Canonical Core IR]
    The single authoritative representation of a program in Lumen.
    A graph of typed nodes, identified by content hashes.
    Acyclic for non-recursive definitions; cycle handling is
    deferred.  In the current phase, canonicalization is structural
    (alpha-equivalence only).

  \item[Projection]
    A rendering of the Core~IR into a specific format for a specific
    audience.  Text, diagrams, and LLM-optimized serializations are
    all projections.  A projection is \emph{lossless} if the Core~IR
    can be exactly reconstructed from it, and \emph{lossy} otherwise.

  \item[Human Projection]
    A projection optimized for human reading and writing.  Typically
    a textual notation resembling a conventional programming language.

  \item[LLM Projection]
    A projection optimized for consumption by a large language model.
    Its design is an open research question; the 90-day experiment
    compares candidate formats.

  \item[Content Addressing]
    Identifying code units by the cryptographic hash of their
    canonical form rather than by name or file path.  Names are
    aliases for hashes.  Adopted from the Unison language.

  \item[Structural Canonicalization]
    Reducing a program representation to a normal form in which
    alpha-equivalent programs are identical.  Achieved by converting
    bound variables to de~Bruijn indices.

  \item[Semantic Canonicalization]
    The (unachieved, long-term) goal of reducing semantically
    equivalent programs to identical normal forms.  Undecidable in
    general.

  \item[Staged Strictness]
    The principle that a Lumen program can range from fully untyped
    to fully dependently typed, with the level of strictness chosen
    per module.  An architectural default of the project.

  \item[Contract]
    A formal specification of a function's expected behavior,
    expressed as preconditions, postconditions, and invariants.
    In the current experimental phase, contracts are generated by an
    LLM and reviewed by the author per the contract review protocol.

  \item[Semantic Graph]
    Synonym for Core~IR when emphasizing its graph structure.

  \item[Semantic Diff]
    A diff algorithm operating on Core~IR graphs rather than text
    lines.  Resistant to renames and formatting changes.  Deferred.

  \item[Lossless Projection]
    A projection from which the original Core~IR can be exactly
    reconstructed.  At least one must exist at all times.

  \item[Route A]
    The empirical research path: experiments comparing code
    representations for LLM tasks.  The 90-day study; completed in v0.4.

  \item[Route A2]
    The continuation of the empirical path after Route~A: ceiling-free
    bounded-regime testing plus a Semantic Coding Layer prototype.  The
    active phase as of v0.4.

  \item[Route B]
    The verification-core research path: formal design of the
    Core~IR with dependent types, contracts, and verification
    infrastructure.  Deferred to a future phase.

  \item[Architectural Default]
    Term used in v0.2.1; superseded by ``Long-Term Research Bet'' in
    v0.2.2.

  \item[Long-Term Research Bet]
    A design principle that the project adopts as a strong directional
    commitment.  It is not an empirical claim and is therefore not
    directly falsifiable by a single experiment.  However, it must
    declare an evidence threshold at which the project will formally
    reassess the bet and its motivating rationale.  Replaces
    ``architectural default'' from v0.2.1.

  \item[Information-Parity Control]
    The experimental condition C1\textsuperscript{+}, which carries
    the same information as C4 (types and contracts) but in textual
    rather than structured form.  Required to separate the effect of
    representation \emph{structure} from the effect of information
    \emph{enrichment}.

  \item[Primary Endpoint]
    The single most important test in the confirmatory analysis
    family: H1 tested on T2 (bug detection composite score).  Other
    confirmatory tests are also reported but carry less interpretive
    weight.

  \item[Confirmatory Analysis]
    Pre-registered statistical tests (H1--H3) with family-wise error
    correction (Holm--Bonferroni).  Distinguished from exploratory
    analysis.

  \item[Exploratory Analysis]
    Post-hoc or non-pre-registered statistical analyses, reported
    with uncorrected $p$-values and explicitly labeled as exploratory.
    No causal or confirmatory claims are made from these results.
\end{description}

\end{document}
